% =====================================================================
%  05-identidad.tex - Verificacion de identidad, sesiones y cuentas
% =====================================================================

\section{Verificación de identidad, sesiones y cuentas}\label{sec:pruebas-identidad}

Esta sección reúne la verificación de la autenticación y del ciclo de vida de las cuentas descritos en la \cref{sec:impl-identidad}: diagnóstico del segundo factor, directorio de miembros, plazos de sesión y recuperación de borradores de Qadra tras un nuevo acceso, recuperación de contraseña, invitaciones, vínculo y acceso por certificado del Owner y calibración del costo de las credenciales.

El registro operativo de los intentos de segundo factor (\cref{sec:impl-s8-totp}),
común al acceso por contraseña y por certificado, tiene una verificación focal
propia. La del 7 de octubre de 2026 aprobó 34
pruebas de identidad, 14 HTTP y 22 del acceso por certificado, 70 en total, en
Windows, con adaptadores de memoria o sustitutos HTTP. Las del acceso por
certificado conservaron las comprobaciones de autoridad y admisión de sesión. La campaña incluyó
aceptación, código inválido, reutilización, desafío ausente, cuenta inactiva,
cambio de generación durante la verificación, recuperación y fallo de auditoría.
La captura HTTP comprobó correlación por petición, rechazo público genérico y
ausencia de secretos en los eventos. Este ensayo no acredita persistencia real,
despliegue ni una nueva regresión global; su alcance está registrado en
\rutarepo{docs/verification-report.md}.

\subsection{Directorio de miembros y revocación durable}\label{sec:pruebas-miembros}

Las campañas de esta subsección verifican el directorio de miembros y la revocación durable de la \cref{sec:impl-miembros} y corresponden a la verificación del 19 de septiembre de 2026. La campaña focal de aplicación y HTTP aprobó \textbf{41 pruebas en 42.282~s}: veinte de identidad, cuatro de consultas, diez de servicio y siete de transporte HTTP. Comprueba generaciones obsoletas, filtros y cursores ligados al contexto, denegación a los otros tres roles, reautenticación del principal completo, proyecciones inválidas, conflictos y respuesta confirmada de un cambio propio. Los fallos de auditoría no se presentan como éxito.

La campaña de PostgreSQL y Redis aprobó \textbf{12 pruebas en 41.896~s}, con 2553 fuentes estables. Incluye filtrado previo a paginación, asignaciones inactivas conservadas, avance único de generación, protección del último Owner y dos degradaciones concurrentes. También comprueba rollback conjunto ante fallo de auditoría, rechazo de escrituras directas inválidas, catálogo y privilegios, inventario sin Owner activo, valores Redis anteriores o fuera de rango y restauración de generaciones, asignaciones y guardas mediante respaldo real. Esta restauración focal no sustituye la aceptación HTTP completa después de restaurar.

El cliente aprobó \textbf{16 pruebas Node en 0.655~s}; el navegador con HTTP controlado aprobó \textbf{11 escenarios en 28.848~s}. Cubren directorio y filtros, revisiones decimales sin pérdida de precisión, respuestas inciertas, selección por expediente, permisos, último Owner y limpieza tras cambio propio o respuestas tardías. El alta directa es distinta de una invitación. Estos recorridos no constituyen evidencia de persistencia real ni una evaluación con usuarios.

El navegador con servicios reales aprobó \textbf{tres recorridos en 418.426~s de comando y 40.1~s de Playwright}, con 2553 fuentes estables. Comprobó la operación a 1440 y 390 píxeles, incluida la administración de miembros en un expediente cerrado, el rechazo del token anterior después de desactivar y reactivar la cuenta y las restricciones por rol, incluido Client. Las cuatro capturas de escritorio y móvil fueron inspeccionadas: resultaron legibles, sin desbordamientos ni solapamientos y conservaron el lenguaje visual de Qadra. La inspección no constituye una evaluación de usabilidad con personas.

La aceptación API completa con restauración aprobó en \textbf{453.174~s}, con 2553 fuentes estables. El respaldo y la restauración finales conservaron las proyecciones de cuentas y asignaciones; los tokens anteriores recibieron 401 y el nuevo acceso exigió MFA. La comparación mantuvo idénticos los ZIP de cuatro documentos importados y sus 70 eventos históricos de auditoría. El inventario final comprendió 19 raíces documentales, 25 versiones y un incidente de integridad.

En la revisión de esas campañas, los controles de formato, lint y versión mínima de Rust aprobaron en CI, sin sustituir las comprobaciones globales; en ese corte quedaban pendientes la regresión global de CI y la compilación e inspección del PDF que incorporara estos resultados, cada una con su propia comprobación. La integración del incremento en \texttt{main}, con su cierre global, consta en \rutarepo{docs/product-completion.md}; las cifras anteriores describen la aceptación original y no una nueva ejecución. Los registros por campaña se conservan en \rutarepo{docs/verification-report.md}; no se atribuye una nueva suite global a estos resultados. La recuperación de contraseña, las invitaciones y el acceso por certificado tienen su propia evidencia (\cref{sec:pruebas-recuperacion-password,sec:pruebas-invitaciones,sec:pruebas-vinculo-certificado}). No existe un recorrido para restablecer o volver a enrolar un segundo factor perdido; los códigos de recuperación de la \cref{sec:impl-http-identidad} sólo completan accesos individuales.


\subsection{Plazos de sesión, reingreso y recuperación de borradores}\label{sec:pruebas-sesiones}

La aceptación del 2 de octubre de 2026 distingue la política predeterminada
de veinticuatro horas del límite opcional de inactividad. Las pruebas iniciales
reprodujeron rutas ausentes, aceptación de claves antiguas o sin expiración,
metadatos MFA incompletos y configuración desconocida. Después de implementar
el contrato aprobaron 31 pruebas de aplicación, 13 del contrato HTTP y cuatro
del parser real de configuración. Estos grupos utilizan dobles cuando cruzan
los puertos de identidad y persistencia.

Otras 23 pruebas ejercitaron el adaptador contra Valkey 8.1.10 desechable,
autenticado y restringido a retorno local. Comprobaron formato, política,
expiración exacta, precisión de generaciones, revocación concurrente y ausencia
de renovación por lectura. El uso de \texttt{PEXPIRETIME} requiere Redis 7 o
posterior; esta ejecución local no se presenta como una medición de Redis 7.4.

La demostración API integrada aprobó en 320.859~s, incluida compilación,
preparación y restauración, con la política predeterminada. Una aceptación
focal posterior habilitó cinco segundos de inactividad sólo en una instalación
aislada con PostgreSQL, Valkey y material criptográfico real. Aprobó
139 comprobaciones sobre 29 peticiones en 18.306~s: MFA mediante TOTP y código
de recuperación, metadatos públicos, lecturas sin renovación, rechazo de cuerpo
y consulta, extensión explícita hasta el límite absoluto, expiración y logout
sin recreación de claves. Se retiraron los procesos y archivos temporales
propios. Los cinco segundos son un parámetro de prueba, no una decisión de
operación del despacho.

El cliente HTTP, reloj y monitor independientes aprobaron 53 pruebas.
Dos casos reprodujeron primero la recuperación indebida de tiempo descontado
por latencia cuando una respuesta posterior era más rápida. El monitor
corregido mantiene la conversión temporal más conservadora durante la sesión;
no genera actividad autónoma, consultas periódicas ni reintentos. La campaña
remota del contrato de backend aprobó 3381 pruebas Rust, 419 de navegador con
respuestas simuladas y 51 contra servicios reales, además de los umbrales de
cobertura de 97/95/93\,\%. Esa campaña precede a los adaptadores de recuperación
de formularios y no acredita su comportamiento.

El resto de esta subsección reúne, por familia de editores, la verificación del
reingreso de Qadra y de la recuperación de borradores tras un vencimiento o una
nueva autenticación (\cref{sec:impl-qadra-sesion,sec:impl-qadra-borradores});
cada campaña conserva su propio alcance y no se suma a las demás.

La integración posterior del navegador cuenta con 45 comprobaciones focales
del ciclo de sesión, registro de borradores y alcance del descarte por
denegación. Los recorridos simulados verificaron vencimiento sin otra respuesta
HTTP, reautenticación, retirada de autorización, cierre concurrente y comparación
de clasificación. También reprodujeron y corrigieron dos carreras: una navegación
durante la comprobación de sesión destruía el editor; una respuesta de logout
pendiente mantenía bloqueado el control de cierre de una cuenta autenticada
después. La comprobación de visibilidad controla eventos y relojes del navegador;
no representa una suspensión física del sistema operativo.

Tres recorridos focales de expediente aprobaron la conservación de texto parcial
y delitos aún sin confirmar, el descarte al cambiar de cuenta y la comparación
de una revisión administrativa nueva antes de un envío explícito. Seis
recorridos de carga documental principal comprobaron bytes y nombre del archivo,
campos pendientes, cancelación, cambio de cuenta, cierre del expediente y pérdida
de autorización. Una carga con respuesta perdida permaneció bloqueada hasta
consultar el listado actual y aceptar explícitamente el riesgo de duplicación.
Otros cuatro recorridos aprobaron en 16.2~s: el primer montaje después de MFA
oculta esperó la confirmación visible; la nueva versión documental recuperó
archivo y nombre sólo tras consultas frescas, conservó la base original ante
una versión concurrente y descartó la captura al cancelar o cambiar de cuenta.
Estos resultados usan respuestas HTTP simuladas y no sustituyen una campaña
integrada de todas las familias de editores. La recuperación universal y la
activación operativa del límite de inactividad permanecen pendientes.


La ampliación de participantes fue aceptada mediante 25 escenarios distintos
repartidos en campañas focales: ocho de ficha manual, nueve de identidad
representada y ocho de ficha tipificada. Incluyen autorización nueva, cambios de
revisión, cierre, rechazo de acceso, resultados inciertos y retiro de la captura
al confirmar. Ocho escenarios de etapas aprobaron en 28.4~s, con historia completa,
selecciones y archivos por propietario y sin reenvío automático. Estos grupos
usan HTTP controlado y no se presentan como una única regresión global.

Un recorrido adicional con servicios desechables aprobó en 37.5~s: atravesó dos
vencimientos y tres autenticaciones MFA, conservó campos de expediente y el archivo
principal, y exigió acciones explícitas para las dos escrituras confirmadas. La
preparación y ejecución completas tomaron 102.735~s, incluida compilación de
51.8~s. El wrapper se incorpora a la campaña remota con diagnóstico y artefacto
obligatorios. La duración de inactividad usada en ese ensayo no habilita ni fija
la política del despliegue privado; los otros editores mantienen su inventario
de aceptación separado.


Diez escenarios de audiencias y resultados aprobaron en 33.8~s con un trabajador
local. Abarcaron autorización nueva, referencias históricas, comparación de base,
archivos por propietario, sesiones, continuaciones e incertidumbre resuelta sólo
por recibo exacto. Trece escenarios existentes de programación, referencias y
selector compartido aprobaron por separado en 30~s. Se mantuvieron tiempos y
reglas del dominio; la aceptación usa respuestas HTTP controladas y no acredita
por sí sola regresión remota ni recuperación de estas familias con servicios reales.


Ocho escenarios de resoluciones y notificaciones aprobaron en 29.0~s con un
worker y respuestas HTTP controladas. Conservan campos y archivos tras
vencimiento, referencias históricas y base original, y exigen nueva autorización
antes de mostrar el borrador. Incluyen incertidumbre, recibo ajeno rechazado,
limpieza antes del refresco confirmado, referencias transitoriamente inaccesibles
y cambios de estado del expediente. El caso de reapertura reprodujo un estado
cerrado obsoleto en la vista contenedora y aprobó después de sincronizarlo,
sin modificar su aserción ni plazo. Estas pruebas no acreditan todavía el
recorrido de ambas familias contra servicios reales ni su integración remota.


La recuperación de acceso de miembros aprobó seis escenarios nuevos en
18.3~s; once recorridos existentes aprobaron en 25.1~s y dieciséis pruebas
Node de clientes en 0.232~s. Se comprobó autorización Owner y destino antes
de mostrar valores, revisión decimal exacta, resultado incierto sin repetición,
limpieza anterior a la expiración por cambio del propio rol y descarte ante
pérdida de permiso, otra cuenta o cierre explícito. Estos resultados utilizan
HTTP controlado y corresponden al recorrido administrativo del Owner
(\cref{sec:pruebas-miembros}); no
incluyeron el alta por invitaci\'on, cuyas pruebas se presentan por separado
en la \cref{sec:pruebas-invitaciones}.

La recuperación del alta de integrantes aprobó veinte pruebas Node en
298.159~ms: seis nuevas de captura individual y catorce regresiones del
registro. Comprobaron proyección síncrona clonada, exclusión de secretos,
rechazo de registros obsoletos y revalidación de cuenta, identidad y generación
después de proyectar. Capturar una entrada no reemplazó su instantánea pendiente
ni invalidó la autorización de recuperación de otro editor.

Nueve recorridos de navegador aprobaron en 22.5~s con un trabajador y HTTP
controlado: ocho nuevos y el existente de alta MFA y auditoría móvil. Incluyeron
correo parcial y rol sin contraseña, autoridad nueva, descarte por otra cuenta
o logout, rechazo de permiso y limpieza anterior al refresco confirmado. Un
envío incierto no se repitió ante presencia o ausencia del correo exacto en el
directorio paginado. La navegación, antes y después del reingreso, mantuvo la
intención hasta su descarte explícito; las respuestas tardías no recuperaron
el material MFA. Los casos finales incluyen los recorridos previos y no se
suman como identidades adicionales. Esta evidencia local no acredita nueva
aceptación con servicios reales, restauración operativa, integración remota ni
recuperaci\'on del material MFA del alta directa. El alta por invitaci\'on
conserva su propia evidencia en la \cref{sec:pruebas-invitaciones}.

Los calendarios jurisdiccionales aprobaron seis escenarios nuevos en 19.1~s
con respuestas HTTP controladas. La recuperación consulta primero el Owner
y el catálogo o la cabecera vigente; conserva campos parciales, identidad de
fuentes y excepciones, base original y envío incierto. Una revisión ajena no
confirma la operación y un recibo propio retira el borrador antes del refresco.
El contexto es global y no hereda acceso de un expediente seleccionado.
Esta evidencia no valida jurídicamente las fuentes declaradas ni habilita
la política de inactividad del servidor.
Diecinueve recorridos anteriores aprobaron por separado en 38.8~s, incluidos
permisos, conflictos, historia, fuentes y diseño a 390 y 1440 píxeles.

Los recursos procesales aprobaron ocho escenarios nuevos en 27.8~s y doce
regresiones en 26.6~s. Los vínculos de actividades aceptaron ocho escenarios
nuevos en dos focales, con seis de selección y resultado en 21.5~s; doce
regresiones adicionales aprobaron en 21.8 y 11.2~s. Se comprobaron referencias
históricas, permisos nuevos, archivos separados por propietario donde existen,
selecciones parciales sin aprobación y consulta del recibo antes de decidir
ante una respuesta perdida. El borrador confirmado no reaparece después de
otra autenticación mientras el listado sigue pendiente. Las pruebas fueron
locales, con un trabajador y respuestas HTTP controladas; no representan una
campaña integrada ni una aceptación con personas o servicios reales de todas
esas familias.

La creación compuesta de plazos desde recursos aprobó ocho escenarios nuevos
en 34.9~s y trece regresiones en 34.0~s. Se conservaron campos parciales,
identificadores, capturas y base; no se recuperó una aprobación anterior.
La ausencia, un recibo ajeno y una pareja creada por separado conservaron la
incertidumbre. Sólo la confirmación explícita del origen conjunto retiró el
borrador antes de refrescar. La aceptación es local con HTTP controlado; las
confirmaciones administrativas tienen evidencia separada más adelante.

La recuperación del editor ordinario de plazos aprobó ocho casos nuevos en
27.0~s y siete regresiones en 19.7~s, con un trabajador y HTTP controlado.
Incluyó atención con desplazamiento parcial, cambio de revisión, pérdida de
aprobación, cierre y reapertura, ausencia de resultado, recibo ajeno y retirada
del borrador confirmado antes de refrescar. La corrección conservó límites y
aserciones. Estos recorridos locales no acreditan integración ni el recorrido
ampliado con servicios reales.

Seis escenarios de recuperación de solicitudes de informes aprobaron en 16.5~s
con HTTP controlado y un trabajador. Verificaron cuenta y selector actuales,
fechas incompletas, selección retirada, dos vencimientos durante un envío,
repetición explícita de un identificador durable y limpieza tras confirmación.
Otra cuenta, cancelación y pérdida de permiso descartaron el borrador. La
campaña no envió correo ni repitió la aceptación PDF/CSV con servicios reales.

Once recorridos existentes y tres de resultado recuperado aprobaron en 26.2~s
tras retirar un aviso duplicado de revocación. Se conservaron las descargas,
permisos, paginación, estados durables y descarte de respuestas tardías.

Nueve escenarios distintos aceptaron la recuperación de preferencias personales
con HTTP controlado y un trabajador. Cinco verificaron contexto y campos;
los tres resultados y una denegación aprobaron en un focal de 17 recorridos
en 34.8 s, junto a trece confirmaciones independientes. Una carrera del fixture
se corrigió esperando los campos recuperados antes de retener otra consulta.
La denegación de la bandeja descartó la captura todavía sin abrir. Dos regresiones
anteriores de preferencias y paginación aprobaron dentro de otro focal de
21 recorridos en 39.3 s. Estas mediciones no acreditan correo externo ni
activación operativa de inactividad.

Trece escenarios de confirmaciones reprodujeron primero la posibilidad de repetir
estados o sellados tras una respuesta perdida, incluso al cancelar y reabrir.
Después aprobaron en un focal de 17 casos en 34.8 s junto a cuatro preferencias.
Diecinueve regresiones relacionadas aprobaron dentro de 21 casos en 39.3 s.
Verificaron estado vigente, versión exacta, bloqueo durante refrescos, ausencia
de reenvío y conservación del resultado confirmado si una lectura posterior
falla. Son pruebas locales con HTTP controlado y un trabajador; no representan
una nueva campaña integrada ni evaluación con personas.

Otros cuatro recorridos aprobaron en 14.0 s después de vencer durante una
confirmación de estado, participante, asignación o sellado. La nueva MFA precedió
a las lecturas actuales; una sola escritura original quedó registrada y su
respuesta tardía no reemplazó el contexto o la intención nuevos. Las confirmaciones
sin campos editables se descartaron. El reloj del harness se avanzó 20 ms tras
observar la respuesta, conservando los límites del producto y del corredor.

Cuatro escenarios de lectura de avisos aprobaron en 13.8~s con HTTP controlado,
sin modificar el producto. Alertas e informes aplicados o no aplicados exigieron
consultas con la sesión nueva, sin repetir el acuse anterior. Cambiar de cuenta
retiró la tarjeta previa; el informe requirió detalle exacto autorizado. Un éxito
tardío no recuperó un informe denegado y un fallo antiguo no terminó otra
intención pendiente. Estos recorridos conservan separada la aceptación con
servicios reales y no habilitan inactividad operativa.


\subsection{Recuperación interna y generación de autenticación}\label{sec:pruebas-recuperacion-password}

Las pruebas focales del 2 de octubre de 2026 separan el contrato de aplicación
de la persistencia y de los adaptadores criptográficos. Veintiuna pruebas de
aplicación aprobaron con dobles explícitos, incluida la cancelación por ID y
digest exactos para impedir que una respuesta tardía retire otra emisión.
Otra prueba de dominio comprobó la identidad del nuevo tipo de valor.

La primera campaña con PostgreSQL desechable aprobó 19 pruebas en 48.11~s.
Comprobó cupo concurrente, consumo único, rechazo de repeticiones, actividad y
generación vigentes, expiración después de bloqueos y rollback ante errores de
auditoría. Los casos de autoridad verificaron denegación de DML al rol de
ejecución, imposibilidad de confirmar un consumo sin su recibo auditado y
rechazo de cambios en cuerpos, privilegios o catálogo. Otras dos pruebas
aprobaron en 5.07~s: vencimiento durante la espera de la fila de usuario y
resolución del namespace de la tabla validada cuando precede un esquema vacío
sin permiso de creación. Este último caso reprodujo primero una invocación
incorrecta en el primer esquema de búsqueda.

Un recorrido adicional aprobó en 11.79~s con Argon2id, TOTP, protección AES,
PostgreSQL y Redis reales. Tras cambiar la contraseña rechazó el password
anterior, una sesión todavía presente en Redis y desafíos MFA anteriores sin
consumir otro código de recuperación. La contraseña nueva permitió completar
TOTP y recuperación MFA existentes. Se conservaron rol, asignaciones y los
valores MFA que el restablecimiento no debe alterar. Entrega, fuente del token
y limitador fueron dobles declarados; no se envió correo externo.

La invalidación administrativa después de restaurar aprobó doce pruebas en
46.46~s. Un dump y su restauración reales conservaron credenciales y recibos
consumidos, mientras las capacidades pendientes quedaron canceladas junto con un
único evento. Se verificaron reintentos con el predecesor original, rechazo de
alteraciones, rollback y lectura del reloj después de bloqueos. Una prueba
reprodujo una función de bloqueo sombreada en el prefijo de búsqueda; fijar el
catálogo seguro antes de resolver relaciones eliminó esa invocación.

Esas campañas aceptaron la frontera interna. Por sí solas no demuestran
recuperación por HTTP, correo operativo, aceptación humana, política de
restauración completa ni despliegue. Las campañas focales no se suman como una
regresión única del workspace; el cierre integrado conserva sus propios gates.


La frontera pública dispone de campañas focales separadas: catorce pruebas de
correo con proveedor simulado, diez de HTTP aislado y cinco de composición con
presupuesto compartido. La respuesta del proveedor debe ser un objeto JSON;
una prueba reprodujo la aceptación incorrecta de arrays posicionales antes de
corregirla. Once recorridos distintos de navegador, repartidos en cuatro
campañas, comprobaron retiro del fragmento, ausencia de envío automático,
confirmación de contraseña, incertidumbre y descarte de respuestas tardías.
Cinco pruebas del enlace complementaron esa verificación.

El consumidor, la configuración y su supervisor aprobaron diecinueve pruebas;
dos adicionales comprobaron bind fallido, cierre conjunto y vida de los
propietarios síncronos fuera del runtime. Seis pruebas del binario verificaron
las opciones y la independencia de canales de correo. La demostración existente
de CLI y la API integrada aprobaron con recuperación deshabilitada; la última
conservó cuatro documentos, setenta eventos históricos y el ZIP de evidencia
idéntico tras restaurar. Este resultado protege los recorridos anteriores, pero
no sustituye la aceptación HTTP específica de recuperación.

La composición HTTP de recuperación aprobó después un recorrido en 13.31~s,
con compilación de 10.53~s. Usó PostgreSQL, Redis, entropía del sistema, Argon2id,
TOTP y códigos MFA reales; sólo se sustituyó la entrega por una captura local.
Verificó la retención del presupuesto compartido, cambio auditado, consumo único
e invalidación de contraseña, sesiones y desafíos anteriores. La contraseña
nueva exigió el segundo factor existente y no alteró rol ni pertenencias.
No se envió correo externo ni se activó la función en el despliegue.

Los comandos administrativos de restauración aprobaron nueve casos con el
binario compilado y PostgreSQL real en 19.21~s. Comprobaron argumentos,
predecesor y recibo, reintento, rechazo al rol de ejecución y catálogo sin
reparación mediante conexión de sólo lectura. Otra prueba Unix reprodujo la
exposición de un valor de conexión no UTF-8; un diagnóstico fijo la corrigió y
aprobó en 0.01~s. Estos comandos no coordinan por sí mismos una restauración operacional.

La regresión completa de la frontera interna aprobó 3457 pruebas Rust,
447 de navegador controlado y 51 con servicios reales, además de dos ignoradas
históricas y los gates de cobertura de 97/95/93\,\%. Su confirmación natural
aprobó CI en 10~min~26~s, Web en 12~min~02~s y el documento en 1~min~15~s.
Los adaptadores públicos descritos conservan su evidencia focal independiente:
no se atribuyen a esa regresión, a correo externo ni a una activación en VPS3.

La aceptación productiva se registra por separado en
\rutarepo{docs/production-acceptance.md}. La instalación \texttt{v0.1.6}
conservaba la recuperación pública activada con Resend y los valores aprobados
de la \cref{sec:impl-recuperacion-password}. El 12 de octubre de 2026, el
usuario informó que otra persona había probado el día anterior la recuperación
con su cuenta. Ese reporte humano de uso real no constituye una nueva
campaña instrumental: no se le atribuye la comprobación de cada control
negativo, que conserva la evidencia técnica anterior.


\subsection{Comprobaciones focales de invitaciones}\label{sec:pruebas-invitaciones}

Las comprobaciones terminadas separan dobles de aplicaci\'on, adaptadores nativos
y transporte controlado. Nueve pruebas del protector AES-GCM comprobaron
separaci\'on de finalidad, identidad y revisi\'on, alteraci\'on del texto cifrado,
l\'imites y errores sin secretos. La suite de aplicaci\'on aprob\'o 32 casos:
incluye los 23 del ciclo de invitaci\'on y nueve del consumidor de entrega; esos
grupos no se suman otra vez. Los dobles comprueban autoridad Owner, ausencia de
cuenta antes de TOTP, recuperaci\'on del mismo material MFA, recibo tras respuesta
perdida y rechazo de una intenci\'on distinta.

El consumidor se comprob\'o con un solo intento por ejecuci\'on. Los casos
validaron finalidad y contexto del descifrado, reloj antes del env\'io, expiraci\'on
durante el trabajo criptogr\'afico, estados aceptado, rechazado e incierto,
reintento dentro de la ventana original y resultado obsoleto al perder el lease.
Cinco pruebas TCP locales del correo comprobaron cuerpo e identidad estables,
fragmento privado, respuestas del proveedor, tiempo acotado y ausencia de
redirecciones o reintentos del cliente HTTP. No enviaron correo externo.

Siete pruebas de PostgreSQL aprobaron conservaci\'on del material pendiente,
cuenta y recibo \'unico al aceptar, reenv\'io que invalida trabajadores anteriores,
rechazo de cuentas existentes o Owner revocado y reversi\'on ante fallo de
auditor\'ia, confirmaci\'on concurrente y capacidad vencida. Otras siete
verificaron permisos de columnas, material cifrado
obligatorio y lease limitado por la ventana del proveedor. Dos pruebas Redis
comprobaron presupuesto at\'omico e independiente y errores sin datos privados.

Nueve casos del contrato HTTP y dos de composici\'on aprobaron la forma exacta
de los comandos y respuestas, revisiones decimales, campos desconocidos,
par\'ametros ambiguos, errores sanitizados y l\'imites compartidos. Una petici\'on
cancelada conserv\'o el presupuesto del trabajo que segu\'ia en ejecuci\'on. Las
rutas p\'ublicas devolvieron preparaci\'on o recibo sin emitir una sesi\'on.

Dos recorridos de Owner con respuestas HTTP controladas comprobaron emisi\'on,
reenv\'io expreso y conciliaci\'on de una respuesta perdida tras navegar. Esa
ejecuci\'on fall\'o en un tercer caso, que reprodujo la duplicaci\'on de una
invitaci\'on creada entre dos p\'aginas. Tras corregir la uni\'on por identidad,
el caso de paginaci\'on aprob\'o por separado con cada fila una sola vez,
continuaci\'on exacta y sin errores de p\'agina. Estos resultados no se presentan
como una sola campa\~na de navegador completamente aprobada.

La invalidaci\'on administrativa aprob\'o un caso sobre PostgreSQL: retir\'o dos
invitaciones pendientes, una con preparaci\'on MFA, y conserv\'o cuenta y recibo
aceptados. Rechaz\'o el uso posterior de los enlaces y del lease retirado. La
misma ejecuci\'on aprob\'o doce regresiones de restauraci\'on de contrase\~na.
Otro caso posterior aprob\'o la reversi\'on completa de la invalidaci\'on cuando
se impide insertar su auditor\'ia. Estos resultados comprueban la primitiva
administrativa sobre un estado preparado y se distinguen del ensayo de copia
restaurada descrito a continuaci\'on.
Otros dos casos de guardas comprobaron reintento con recibo original y rechazo
de autoridad, destino o antecedente incorrectos en una ejecuci\'on previa
parcial, cuyo fallo estaba en el fixture de la auditor\'ia. La correcci\'on de
ese fixture y su resultado posterior se mantienen separados.

El ensayo de respaldo poblado aprob\'o despu\'es en 23.014~s, incluidos 6.20~s
del caso. Gener\'o un dump y lo restaur\'o en una base distinta. Compar\'o cuentas,
operaciones y recibos, anul\'o la invitaci\'on pendiente y recuper\'o la aceptada
sin duplicar su cuenta. Conserv\'o el prefijo de auditor\'ia y comprob\'o
criptogr\'aficamente la cadena con el nuevo evento de invalidaci\'on. Los enlaces
y entregas pendientes no volvieron a admitirse; la base de origen permaneci\'o
intacta. Este ensayo complementa las pruebas unitarias y de puertos con una
transferencia real de estado entre bases.

La composici\'on del binario aprob\'o despu\'es cinco casos en 27.48~s, con
58.1~s de compilaci\'on y ejecuci\'on. Uno recorri\'o el router HTTP con
PostgreSQL y Redis reales: preparaci\'on sin cuenta, confirmaci\'on TOTP,
repetici\'on y consulta del mismo recibo, reapertura, reenv\'io con rechazo del
enlace anterior y acceso normal con MFA. La cuenta aceptada permaneci\'o \'unica
y la cadena auditada conserv\'o sus eventos esperados. La entrega se captur\'o
mediante un doble, sin enviar correo externo. Los otros cuatro casos cubrieron
activaci\'on independiente, pausas del consumidor, cierre supervisado y argumentos
administrativos de restauraci\'on.

Las demostraciones completas de API y CLI del repositorio aprobaron en
594.352~s y 28.901~s, respectivamente. La regresi\'on de navegaci\'on de
recuperaci\'on de contrase\~na aprob\'o en 8.662~s; el lint final y la
compilaci\'on web aprobaron en 8.08~s y 4.225~s. Se registran como verificaciones
separadas y no se suman a los casos focales de invitaciones.

Los resultados focales se registran por separado en
\rutarepo{docs/verification-report.md}. El resultado funcional local confirma
alta con MFA, reentrada mediante recibo y conservaci\'on de cuentas y evidencia
tras restaurar, sin duplicar aceptaciones. Ninguna de estas comprobaciones
envi\'o correo externo, cuya entrega permanece deshabilitada por defecto. Las
pruebas del router, del navegador controlado y de restauraci\'on conservan su
alcance propio y no acreditan CI, integraci\'on ni despliegue. La integraci\'on
de las invitaciones con sus gates, su despliegue en \texttt{v0.1.4} y la
activaci\'on de Resend, cuya recepci\'on de un \'unico correo de prueba confirma
el transporte sin acreditar un alta real en producci\'on, se registran en
\rutarepo{docs/product-completion.md}.


\subsection{Registro canónico y verificación del certificado Owner}\label{sec:pruebas-vinculo-certificado}

Las campañas de esta subsección recorren por capas el vínculo y el primer factor
de la \cref{sec:impl-vinculo-certificado}: dominio y verificador, aplicación,
persistencia, transporte y composición, aceptación HTTP conjunta e interfaz, y
después las mismas fronteras del primer factor por certificado.

Ocho pruebas de dominio comprobaron vectores literales de 150 bytes, propósito,
cuenta, contadores y retiro terminal con preservación del registro. El verificador
criptográfico añadió ocho pruebas que aprobaron en 1.00~s, después de 19.02~s de
compilación. Se generaron certificados y firmas con OpenSSL para comprobar
aceptación del perfil exacto y rechazo de otra cuenta, revisión o propósito,
material alterado, revocación y límites inclusivos de vigencia. La confianza
recalculada tuvo que coincidir por completo con la instantánea suministrada.

Las 24 pruebas existentes de declaraciones personales aprobaron en 1.15~s y
conservaron el rechazo de usos extendidos de clave. La comprobación estática
focal aprobó con advertencias tratadas como errores. Estos resultados locales
no prueban la procedencia de una instantánea construida por un llamador,
transacciones auditadas, una asociación persistida ni autenticación por
certificado. La comprobación remota completa se registra por separado.

La capa de aplicación añadió catorce pruebas con puertos controlados, aprobadas
en 0.00~s después de 9.58~s de compilación. Verificaron identidad actual,
reautenticación del principal completo, tiempo posterior a la comprobación
criptográfica, evidencia exacta, UUID solicitado y conservación del primer
retiro concurrente. Una regresión adicional ejercitó el puerto de aplicación
con el verificador real: nueve casos aprobaron en 1.11~s; las 24 regresiones de
declaraciones aprobaron en 1.27~s. Clippy focal aprobó en 14.54~s. Estos ensayos
no sustituyen la aceptación del adaptador PostgreSQL, su auditoría atómica ni
un recorrido HTTP de autenticación por certificado.

La campaña posterior de persistencia, del 3 de octubre de 2026, aprobó doce de
doce pruebas en 47.39~s, tras 3.68~s de compilación, con PostgreSQL~16.15 privado,
autenticación SCRAM y rol restringido. Ejercitó altas, retiros y renovación,
recibos exactos, carreras entre transacciones, huellas permanentes por cuenta,
rollback ante rechazo de auditoría y relectura de autoridad, confianza y reloj
después de esperar bloqueos. Los ocho casos iniciales habían aprobado en
30.39~s, tras 18.23~s de compilación; la corrida final incorpora cuatro casos
adicionales.

Un ensayo reprodujo contadores actuales inferiores a las capturas históricas;
la regresión de cambio legítimo de rol y actividad ya aprobaba. Otros dos
reprodujeron la admisión indebida del permiso para omitir disparadores mediante
el modo de replicación, directo o alcanzable con otro rol. Las correcciones
rechazan esas configuraciones sin cambiar el historial. La comprobación estática
focal aprobó en 6.751~s. El aviso conocido de compatibilidad futura de Redis
permanece separado de estos resultados.

La identidad y el reloj se controlaron en las pruebas; PostgreSQL, la publicación
de confianza y RSA fueron reales. La admisión vuelve a verificar las firmas con
su confianza e instante capturados y comprueba la correspondencia exacta con
los eventos. Las pruebas verifican además la cadena global mediante su
comprobador existente. La modificación administrativa de una fila simuló un
inventario incoherente; no se ejecutó un volcado y restauración completos. Estos
resultados no acreditan enrolamiento, un ingreso MFA real, rutas HTTP, acceso por
certificado ni firma documental individual.

La aceptación independiente del adaptador HTTP añadió doce pruebas con el
servicio real de aplicación y puertos controlados. Aprobaron en 0.03~s después
de 28.12~s de compilación; comprobaron las cuatro rutas, JSON estricto, límites
y Base64 canónico, autoridad actual, precisión de evidencia, errores neutrales
y respuestas sin almacenamiento en caché. Estas solicitudes se ejecutan contra
el router en proceso, sin RSA ni PostgreSQL reales.

La composición reprodujo primero la ausencia del campo requerido para inyectar
el servicio. La campaña posterior aprobó veinte de veinte casos, tras 22.78~s
de compilación: tres nuevos de composición en 0.33~s, doce del adaptador
independiente en 0.03~s y cinco regresiones de composición de recuperación de
contraseña en 0.34~s. Los nuevos casos retienen trabajo síncrono mediante puertos
controlados para comprobar el mismo presupuesto externo, cancelación HTTP en
ambos sentidos entre rutas Owner y existentes, rechazo antes de leer un cuerpo
y autorización actual dentro del permiso compartido.

El binario aprobó diez pruebas de opciones después de 68~s de compilación:
seis de recuperación de contraseña en 0.17~s y cuatro de sesión en 0.00~s.
Clippy de composición y binario aprobó en 26.34~s con advertencias como errores;
permaneció el aviso conocido de compatibilidad futura de Redis. Las pruebas de
opciones no iniciaron el servidor.

La composición del binario utiliza los adaptadores reales, pero esa campaña
no inicia el proceso servidor ni reúne autenticación MFA, PostgreSQL y RSA en
un recorrido HTTP. La evidencia de persistencia y criptografía anterior conserva
sus propios controles y límites; los resultados no se suman como si probaran
esa integración ni una instalación operativa. El acceso por certificado y la
firma documental individual permanecen fuera de estos ensayos.

La aceptación HTTP conjunta del 3 de octubre de 2026 concluyó con salida cero
en 334.419~s mediante \rutarepo{scripts/api-demo.sh}, incluida la compilación
y la preparación del entorno desechable. Usó PostgreSQL~16.15, Redis, ingreso
por contraseña y MFA reales, RSA y la CA interna. Reutilizó qpdf~12.4.1 y los
decodificadores multimedia previamente verificados; no se registró la versión
del servidor Redis. Esta campaña conserva por separado las mediciones focales.

Un titular sintético firmó fuera de HTTP los 150 bytes canónicos, reconstruidos
independientemente, y se verificó con OpenSSL la firma pública devuelta; el
ensayo comparó esa prueba con las filas y los eventos auditados. La firma
alterada produjo rechazo sin escritura; un alta válida añadió un único evento. Las consultas y repeticiones
exactas conservaron el recibo y los contadores de cuenta. Una CRL sucesora
revocó la hoja de ensayo: la preparación anterior produjo conflicto y una
nueva firma con el certificado revocado fue rechazada. El recibo original
siguió disponible y su retiro conservó la evidencia sin exigir otra firma.

El volcado y la restauración SQL preservaron ambas tablas y sus enlaces de
auditoría. Tras invalidar sesiones y desafíos anteriores y renovar MFA,
consulta y repetición devolvieron el mismo recibo terminal, con declaración,
certificado, firma, confianza y fechas originales, sin eventos adicionales ni
reactivación. Las cuatro rutas rechazaron el token anterior y al Paralegal
vigente sin mutar evidencia. La campaña completa conservó también los
recorridos previos; su comprobación final mantuvo cuatro documentos y los 70
eventos del prefijo importado, con evidencia ZIP idéntica.

El resultado acredita ese recorrido local con servicios reales, sin publicar ni
instalar la extensión y sin la interfaz Qadra, cuya verificación se presenta en
los apartados siguientes. No acredita inicio de sesión por certificado, firma
individual de documentos, custodia exclusiva del titular, prestación de un PSC ni restauración operacional
conjunta de SQL, RDB y PKI.

La consulta posterior del vínculo propio sin retirar reprodujo primero la
ausencia del puerto y del servicio. Después aprobaron 27 de 27 pruebas de
aplicación en 0.01~s y 15 de 15 HTTP en 0.04~s con puertos controlados. Dos
casos PostgreSQL reales aprobaron en 6.13~s, tras 13.03~s de compilación; el
ejecutor completo duró 20.472~s y confirmó la retirada del clúster desechable.
Clippy focal aprobó en 24.01~s. Los nueve casos nuevos se distribuyen en
cuatro de aplicación, tres HTTP y dos de persistencia; no se suman otra vez
a los totales que ya los incluyen.

Se comprobaron el recibo propio o su ausencia, la reautenticación del principal
completo en ambos resultados y el rechazo de evidencia ajena o retirada.
PostgreSQL conservó historia, filas y auditoría durante consultas, retiro y
renovación; la confianza vencida no ocultó el vínculo sin retirar. Al desactivar
la cuenta durante una espera real del bloqueo de auditoría, la consulta rechazó
la autoridad posterior, con vínculo y sin él. Esta quinta ruta no formó parte
de la campaña conjunta de 334.419~s: los focales HTTP y PostgreSQL mantienen
sus fronteras separadas y no acreditan otra aceptación integrada ni interfaz.

\subsubsection{Cliente y navegador con transporte controlado}

Los diez casos del cliente aprobaron en 502.558~ms después de reproducir la
ausencia de la fachada y de los helpers de material público. El navegador
reprodujo primero la ausencia de la navegación personal. Tras implementar la
interfaz aprobaron seis recorridos en 16.8~s, con un solo worker y respuestas
HTTP controladas. Registro, retiro, restricciones por cuenta, recuperación de
la intención y consulta de una escritura incierta conservaron UUID y evidencia
exactos. El caso añadido durante revisión rechazó una respuesta terminal cuya
firma original no coincidía; no se acredita un RED independiente para ese caso.

El ensayo de reingreso comprobó primero la identidad vigente y después el
vínculo actual, sin recuperar el editor antes de ambas respuestas. La ausencia
de recibo permitió únicamente un reenvío explícito del mismo comando; el
resultado confirmado después de otra MFA cerró la intención sin otra escritura.
Cerrar sesión o cambiar de Owner descartó la captura. Estos dobles de transporte
no comprueban RSA ni persistencia reales. Dos casos del planificador aprobaron
en 143.671~ms y conservaron las asignaciones previas al incluir un recorrido
real nuevo; la selección por sí sola no demuestra su ejecución.

\subsubsection{Recorrido Qadra con servicios reales}

La aceptación focal posterior aprobó un recorrido de navegador en 12.1~s de
Playwright. El comando completo duró 252.134~s, con 38.39~s de compilación y
la preparación de las familias de la partición seleccionada; el filtro limitó
el navegador, no eliminó esos preparativos. Hubo un compilador y un worker.
PostgreSQL~16.15 y Valkey~8.1.10 ejecutaron persistencia e identidad; este último
fue el almacén compatible con Redis invocado por la herramienta de ensayo. Se
reutilizaron qpdf~12.4.1 y los decodificadores verificados. La familia nueva
preparó en 4.77~s un Owner y certificado Partner propios.

La interfaz descargó la declaración canónica, reconstruida independientemente,
y seleccionó la firma producida por OpenSSL fuera del navegador. Registro,
consulta y descarga conservaron la evidencia pública exacta. Después de cerrar
sesión y completar MFA nueva, la consulta actual descubrió el mismo registro
sin repetir preparación ni alta. El retiro conservó ese registro, produjo la
declaración terminal exacta y permitió descargar su historia; volver a la vista
resolvió la ausencia de vínculo sin retirar. Se ejercitó así la quinta ruta
junto a las cuatro operaciones por UUID y los componentes reales de Qadra.

El resultado utiliza servicios desechables y una clave sintética de ensayo.
No modifica el despliegue activo ni demuestra custodia personal exclusiva,
autenticación por certificado o firma documental individual. La campaña no
repitió la restauración SQL ni la prueba de CRL sucesora del ensayo HTTP
anterior; ambas conservan su evidencia y alcance independientes.

\subsubsection{Estructura y verificación de la declaración del primer factor}

Nueve casos de dominio aprobaron tras reproducir la ausencia del módulo. Dos
vectores literales completos fijaron los 182 bytes y se comprobaron límites,
copias independientes y separación de propósitos. El adaptador posterior
aprobó seis pruebas en 1.13~s, con 13.88~s de compilación. Una firma OpenSSL de
un Partner verificó el formato; registro, retiro, nonce o generación alterados,
confianza discordante, revocación firmada y entradas mal formadas se rechazaron.
Estos resultados no demuestran consumo único del desafío, MFA, sesiones derivadas
ni activación de acceso por certificado; las tres primeras fronteras se comprueban
en los apartados siguientes y la activación operativa, en la
\cref{sec:pruebas-portal-ca}.


\subsubsection{Admisión interna, autoridad y capacidades temporales}

La aplicación aprobó inicialmente 20 casos nuevos y 31 del recorrido de identidad
existente. Dos regresiones posteriores reprodujeron vencimiento absoluto y de
inactividad durante una consulta final de autoridad dentro del mismo segundo.
Tras conservar milisegundos en la admisión, el cierre aprobó los 22 casos nuevos
y los 31 previos, en 0.01 y 0.00~s, tras 4.71~s de compilación.
El puerto criptográfico conservó la evidencia y neutralizó fallos internos en
ocho pruebas, incluidas las seis previas, en 0.77~s. La autoridad PostgreSQL
16.15 desechable aprobó seis casos en 13.88~s: observa cambios de cuenta y CRL
después de esperar el bloqueo auditado, distingue ausencia de fallo SQL y
conserva los recibos históricos. La preparación del cambio de generación se
corrigió para usar revocación y reactivación legales; no se alteró el guard de
base de datos.

Con Valkey 8.1.10 aislado y autenticado, compatible con Redis, aprobaron once
casos de procedencia y sesiones en 0.12~s, 23 regresiones de sesiones de
contraseña en 0.07~s y 21 casos de captura y presupuestos en 2.71~s. Se verificó
consumo con un ganador, vencimiento exclusivo, enteros exactos, límites
compartidos, permisos fallidos sin escritura parcial y rechazo de DTO corruptos.
La primera comparación de no escritura usaba serialización RDB; el diagnóstico
mostró campos y expiración idénticos con orden binario distinto. La comprobación
corregida compara tipo, todos los bytes de campos y plazo absoluto.

Estas campañas focales se ejecutaron en serie con un compilador y un trabajador.
No son una campaña HTTP integral, una restauración autenticada ni una activación
en el servidor. El detalle y la frontera pendiente se conservan en
\rutarepo{docs/owner-certificate-authentication.md}.


La invalidación selectiva de capturas de certificado aprobó diez casos de
comando en 0.016~s y un recorrido nativo en 7.006~s. RDB, AOF y reinicio
conservaron los controles y presupuestos, sin resucitar sesiones, MFA o
capturas retiradas. La lectura de hashes del observador fija RESP2; el primer
intento había recibido el objeto JSON de RESP3 antes del respaldo. El contenido
capturado es público y sintético: esta prueba de persistencia no reemplaza una
restauración completa de autenticación HTTP con RSA y MFA reales.


\subsubsection{Transporte opcional del primer factor}

El transporte HTTP aprobó 19 casos en 0.26~s, tras 33.57~s de compilación;
Clippy focal aprobó en 14.30~s. Se observaron primero los exports y el puerto
faltantes. Las comprobaciones incluyen proyecciones exactas, rechazo de entradas
ambiguas y cuerpos excesivos, errores opacos, compatibilidad de acceso existente
y conservación del presupuesto mientras termina un trabajador cuya respuesta
fue cancelada. La firma admitida sólo entrega un desafío MFA nuevo. El flujo
controlado no acredita RSA, PostgreSQL, Redis, restauración ni activación en
el servidor compuesto; esas fronteras se verifican por separado.


\subsubsection{Aceptación nativa del primer factor compuesto}

Siete casos de configuración y composición aprobaron en 0.09~s y Clippy focal
en 20.99~s. Se comprobó la opción deshabilitada, validación previa y una misma
instancia de identidad para ambos métodos. Un recorrido adicional con
PostgreSQL 16.15 y Valkey 8.1.10 desechables aprobó en 12.63~s, tras 12.42~s
de compilación, con un trabajador. Los servicios privados se retiraron al cerrar.

El ensayo abrió la composición real, registró un Partner y firmó externamente
los 182 bytes. La firma alterada consumió la captura y su repetición fue
rechazada; una firma válida exigió TOTP real antes de crear la sesión. Se
comprobó su procedencia y un plazo independiente de los cinco minutos de la
captura inicial. Una CRL sucesora y el retiro del vínculo rechazaron capturas,
MFA pendientes y sesiones previas de certificado. La contraseña siguió admitida,
las credenciales rechazadas no consumieron recuperación y la auditoría conservó
los recibos originales. Esta prueba usa el router en proceso; el listener HTTP,
la restauración autenticada y la interfaz tienen evidencia independiente, descrita
a continuación.


Tres casos de disponibilidad reprodujeron primero la ausencia de la lectura;
la suite ampliada aprobó los 22 escenarios HTTP en 0.29~s. Se comprobaron
composición presente y ausente, rechazo de cuerpos y consultas, ausencia de
invocación al flujo de identidad y uso del mismo presupuesto de solicitudes.
El recorrido CLI completo también aprobó en 25.997~s, conservando pruebas
positivas y negativas de la evidencia criptográfica y la auditoría.


Un ensayo focal adicional con listener real aprobó en 23.248~s. Conservó los
recibos y auditoría tras restaurar SQL, rechazó tres capacidades antiguas antes
de que expiraran y verificó acceso nuevo con RSA y MFA, junto con contraseña
y MFA independientes. La lectura INFO del observador Redis se corrigió tras
reproducir su interpretación indebida como JSON; su caso nativo aprobó en
0.117~s. Una preparación posterior de restauración rechazó MFA con estado 401.
El diagnóstico controlado reprodujo la reutilización del intervalo TOTP;
tres casos aprobaron en 0.021~s al esperar el intervalo siguiente antes de
crear desafíos, sin borrar marcas de uso ni reintentar MFA rechazada.


La campaña HTTP completa posterior aprobó en 365.645~s con PostgreSQL 16.15,
Valkey 8.1.10, qpdf 12.4.1 y decodificadores multimedia verificados. Conservó
los recorridos de los módulos existentes y sus comprobaciones posteriores a
restauración, además del vínculo Owner y las capacidades del primer factor.
Se comprobaron recibos históricos, controles retenidos, invalidación anterior al
vencimiento y nuevos accesos RSA/MFA y contraseña/MFA. Fue una campaña serial
con un compilador y un trabajador. La interfaz, la regresión CI de la revisión
publicada y la activación operativa conservan fronteras de aceptación separadas.

\subsubsection{Comprobaciones de la interfaz del primer factor}

La verificación del cliente aprobó 36 pruebas Node: once del acceso por
certificado y 25 de intentos de autenticación existentes. Los casos comprobaron
el contenedor público acotado, la descarga exacta de 182 bytes, separación de
propósito, vencimiento y consumo local previo al envío. También exigieron que
una respuesta MFA coherente de otra cuenta o de otro rol no instalara acceso,
y conservaron el descarte de respuestas de intentos anteriores.

Cinco recorridos de navegador con HTTP controlado aprobaron en 20.9~s. Incluyeron
indisponibilidad del servidor, rechazo de contenedores privados, selección de
recibo histórico, firma separada, MFA compartida, respuestas tardías al cambiar
de método y preparación explícita después de un resultado incierto o vencido.
Siete regresiones de autenticación previas también aprobaron. La revisión visual
de cuatro capturas de escritorio y móvil conservó el diseño Qadra sin
desbordamiento observado. Estas comprobaciones no realizan una evaluación de
usabilidad con personas ni sustituyen el recorrido de navegador con RSA,
PostgreSQL y Redis reales o la instalación del servicio.

El recorrido posterior con servicios reales aprobó 1/1 en 14.0~s; el comando
completo, incluida la preparación, duró 208.382~s. PostgreSQL 16.15,
Valkey 8.1.10 y OpenSSL conservaron la firma externa de 182 bytes, MFA obligatoria
y selección de la misma cuenta. El ensayo cerró después esa sesión y volvió
mediante contraseña y otro código de recuperación antes del retiro y consulta
histórica. Conservó las aserciones anteriores. El primer intento había recibido
401 al retirar bajo la autoridad del mismo vínculo; corregir esa preparación
no alteró permisos ni límites. La invalidez por retirada permanece comprobada
por los ensayos propios del backend, sin atribuirla al logout de este recorrido.

Las aceptaciones descritas utilizaron firmas producidas con OpenSSL fuera
del navegador. El ayudante Rust/WebAssembly de
\cref{sec:impl-criptografia-local} conserva esa alternativa y el protocolo
de acceso, pero su ejecuci\'on local y su ciclo de secretos constituyen otra
frontera de prueba. Los resultados anteriores no se atribuyen a ese ayudante ni
demuestran su integraci\'on o activaci\'on operativa.


\subsection{Calibración de Argon2id}\label{sec:pruebas-argon2}

El diseño exige que el costo de un hash nuevo de Argon2id quede entre 0.5 y 1.0 segundos en el hardware de referencia (\cref{sec:autenticacion}). El caso de uso de calibración (subcomando \texttt{auth calibrate}) promedia cinco corridas de hash sobre una muestra fija y clasifica la media contra la banda objetivo inclusiva. En el hardware de referencia (AMD Ryzen 7 7730U, 16 hilos lógicos), la calibración de la demostración de cierre del 17 de agosto de 2026 fijó $m = 256$~MiB, $t = 2$, $p = 1$ y reportó 529.4 milisegundos, es decir, un veredicto dentro de la banda. Esta medición conserva su alcance histórico. La subsección siguiente (\cref{sec:calibracion-credenciales}) reporta, sobre el mismo modelo de equipo pero con otro binario y contexto de ejecución, una medición de dos pasadas por debajo de la banda y la calibración de tres pasadas que se utiliza para hashes nuevos; ninguna de las dos medias de dos pasadas sustituye a la otra.

El ajuste de parámetros está localizado en un único punto del adaptador. Un cambio de hardware o de perfil de binario exige repetir la medición, porque la memoria y el tiempo por hash cambian con el perfil de despliegue, y la banda no describe el rendimiento del acceso completo ni de HTTP bajo carga. El resto del servicio no depende de la banda: el formato PHC con los parámetros registrados, la coincidencia y no coincidencia de contraseñas, el tratamiento del hash corrupto como error y la ausencia de la contraseña en los diagnósticos están fijados por pruebas que pasan en cualquier hardware.

\subsection{Calibración del costo de credenciales}\label{sec:calibracion-credenciales}

El 19 de septiembre de 2026 se repitió la medición del costo de un hash nuevo en un equipo AMD Ryzen 7 7730U, el mismo modelo de la medición histórica de la \cref{sec:pruebas-argon2}, con 16 CPU lógicas y Linux 7.1.13. El binario de desarrollo optimiza las dependencias Argon2 y Blake2 a nivel 3. Las cinco operaciones secuenciales con dos pasadas promediaron 400.3851816 ms, por debajo de la banda de 500 a 1000 ms. Con tres pasadas, el promedio de cinco operaciones fue 659.7596048 ms, dentro de esa banda. Se mantuvieron 256 MiB de memoria, una vía, sal aleatoria de 16 bytes y salida de 32 bytes. Esa base de dos pasadas por debajo de la banda motivó fijar tres pasadas para los hashes nuevos, según \rutarepo{docs/adr/0044-reference-password-hashing-cost.md}.

La comprobación focal aprobó siete casos, incluida la verificación de un PHC histórico de dos pasadas con contraseña correcta e incorrecta. Seis pruebas CLI aprobaron y la demostración completa de comandos terminó satisfactoriamente; su calibración secundaria dio 563.2 ms sobre cinco hashes. Las dos observaciones conservan su contexto. Las identidades del binario, entorno, comandos y resultados se registran en \rutarepo{docs/verification-report.md}.

El ajuste afecta a hashes nuevos; los anteriores conservan sus parámetros y no se reescriben. La banda no mide un inicio de sesión completo ni acredita rendimiento bajo carga. La memoria por operación no aumenta, aunque el trabajo de CPU y la ocupación del trabajador son mayores. Estas comprobaciones focales no constituyen la regresión global; la integración del ajuste en \texttt{main} se registra en \rutarepo{docs/product-completion.md}.

